\begin{CNbox}[unbreakable]{obj}{Learning Objectives}

By the end of this module, you will be able to:

\begin{itemize}
\tightlist
\item
  Explain why 5G abandoned EPC’s point-to-point interface model
\item
  Describe SBA design principles and the role of HTTP/\allowbreak{}2, JSON, and REST
\item
  Identify all 5G Core Network Functions and their services
\item
  Explain the NRF’s critical role in service discovery
\item
  Compare direct vs indirect (SCP) communication models
\item
  Contrast EPC architecture with 5G SBA
\end{itemize}

\end{CNbox}

\Needspace{14\baselineskip}

\Needspace{24\baselineskip}

\CNsection{1}{Why 5G Moved Away from EPC Point-to-Point Interfaces}

\CNchained\CNsubsection{1.1}{The EPC Problem}

In 4G/\allowbreak{}LTE, the Evolved Packet Core (EPC) uses \textbf{point-to-point reference points} between network elements:

{\def\LTcaptype{none} % do not increment counter
\Needspace{13\baselineskip}
\begin{longtable}[c]{@{}cc@{}}
\toprule\noalign{}
\rowcolor{CNink}\CNth{Element Pair} & \CNth{Interface} \\
\CNheadrulesep
\endhead
\bottomrule\noalign{}
\endlastfoot
MME \ensuremath{\leftrightarrow} S-GW & S11 \\
MME \ensuremath{\leftrightarrow} HSS & S6a \\
S-GW \ensuremath{\leftrightarrow} P-GW & S5/\allowbreak{}S8 \\
MME \ensuremath{\leftrightarrow} eNB & S1-MME \\
P-GW \ensuremath{\leftrightarrow} PCRF & Gx \\
\end{longtable}
}

\textbf{The scaling problem:}

\begin{itemize}
\tightlist
\item
  With N network elements, you need up to \textbf{N\ensuremath{\times}(N-1)/\allowbreak{}2} unique reference points
\item
  Each interface has its own protocol specification (Diameter, GTP-C, etc.)
\item
  Adding a new network element requires defining new interfaces to every existing element
\item
  Rigid coupling makes evolution extremely difficult
\item
  Vendor interoperability requires testing every interface combination
\end{itemize}

\textbf{Example:} EPC with 6 core elements \ensuremath{\to} up to 15 unique interfaces, each with its own protocol, message formats, and procedures.

\CNsubsection{1.2}{The SBA Solution}

5G introduces \textbf{Service-Based Architecture} where:

\begin{itemize}
\tightlist
\item
  Network Functions (NFs) expose \textbf{services} via a common interface framework
\item
  Any NF can \textbf{discover and consume} any other NF’s services
\item
  Adding a new NF only requires registering its services — no new interfaces needed
\item
  Common protocol stack (HTTP/\allowbreak{}2 + JSON) for all control plane communication
\end{itemize}

\Needspace{20\baselineskip}

\CNsubsection{1.3}{Cloud-Native Principles}

SBA embraces cloud-native design:

{\def\LTcaptype{none} % do not increment counter
\Needspace{14\baselineskip}
\begin{longtable}[c]{@{}
  >{\raggedright\arraybackslash\hspace{0pt}}m{(\linewidth - 2\tabcolsep) * \real{0.5000}}
  >{\raggedright\arraybackslash\hspace{0pt}}m{(\linewidth - 2\tabcolsep) * \real{0.5000}}@{}}
\toprule\noalign{}
\rowcolor{CNink}\CNth{Principle} & \CNth{Implementation in 5G SBA} \\
\CNheadrulesep
\endhead
\bottomrule\noalign{}
\endlastfoot
\textbf{Microservices} & Each NF is decomposed into independent services \\
\textbf{Stateless} & NF instances don’t store session state locally (externalized to data stores) \\
\textbf{Horizontal scaling} & Add more NF instances behind load balancers \\
\textbf{Containerization} & NFs deployed as containers in Kubernetes \\
\textbf{CI/\allowbreak{}CD} & Independent NF lifecycle management \\
\textbf{Service mesh} & SCP provides routing, load balancing, observability \\
\end{longtable}
}

\Needspace{14\baselineskip}

\CNkeepfig{m18-refpoint-vs-sbi}{8}

\CNsection{2}{SBA Design Principles}

\CNchained\CNsubsection{2.1}{Service-Based Interfaces (SBI) vs Reference Points}

\CNfigure{m18-refpoint-vs-sbi}{EPC reference points versus the 5G service-based bus}

\textbf{Reference Points (4G):} Each line is a unique protocol/\allowbreak{}interface specification\CNbr{}\textbf{Service-Based Interfaces (5G):} All NFs connect to a common bus using the same protocol framework

\CNsubsection{2.2}{NF Service Interaction Models}

NFs communicate using two patterns:

\textbf{1. Request-Response (Synchronous)}

\begin{itemize}
\tightlist
\item
  NF Consumer sends a request to NF Producer
\item
  Producer processes and returns a response
\item
  Example: AMF queries UDM for subscriber data
\end{itemize}

\textbf{2. Subscribe-Notify (Asynchronous)}

\begin{itemize}
\tightlist
\item
  NF Consumer subscribes to events from NF Producer
\item
  Producer sends notifications when events occur
\item
  Example: AMF subscribes to UDM for subscription data changes
\end{itemize}

\Needspace{19\baselineskip}

\CNsubsection{2.3}{HTTP/\allowbreak{}2 as Transport Protocol}

5G SBA chose HTTP/\allowbreak{}2 over HTTP/\allowbreak{}1.1 for critical reasons:

{\def\LTcaptype{none} % do not increment counter
\Needspace{13\baselineskip}
\begin{longtable}[c]{@{}
  >{\raggedright\arraybackslash\hspace{0pt}}m{(\linewidth - 2\tabcolsep) * \real{0.5000}}
  >{\raggedright\arraybackslash\hspace{0pt}}m{(\linewidth - 2\tabcolsep) * \real{0.5000}}@{}}
\toprule\noalign{}
\rowcolor{CNink}\CNth{Feature} & \CNth{Benefit for 5G} \\
\CNheadrulesep
\endhead
\bottomrule\noalign{}
\endlastfoot
\textbf{Multiplexing} & Multiple requests/\allowbreak{}responses over single TCP connection — reduces latency \\
\textbf{Header compression (HPACK)} & Reduces overhead for repetitive signaling headers \\
\textbf{Server push} & Enables efficient notification delivery \\
\textbf{Binary framing} & More efficient parsing than text-based HTTP/\allowbreak{}1.1 \\
\textbf{Stream prioritization} & Critical signaling (e.g., emergency calls) gets priority \\
\end{longtable}
}

\Needspace{17\baselineskip}

\CNsubsection{2.4}{JSON as Data Format}

{\def\LTcaptype{none} % do not increment counter
\Needspace{13\baselineskip}
\begin{longtable}[c]{@{}
  >{\raggedright\arraybackslash\hspace{0pt}}m{(\linewidth - 2\tabcolsep) * \real{0.5000}}
  >{\raggedright\arraybackslash\hspace{0pt}}m{(\linewidth - 2\tabcolsep) * \real{0.5000}}@{}}
\toprule\noalign{}
\rowcolor{CNink}\CNth{Reason} & \CNth{Explanation} \\
\CNheadrulesep
\endhead
\bottomrule\noalign{}
\endlastfoot
\textbf{Human-readable} & Easy debugging, logging, troubleshooting \\
\textbf{Web-native} & Leverages massive web ecosystem tooling \\
\textbf{Tooling} & Validators, editors, parsers available in every language \\
\textbf{Schema support} & OpenAPI/\allowbreak{}JSON Schema for API definition \\
\textbf{Lightweight} & Less verbose than XML (used in older 3GPP specs) \\
\end{longtable}
}

\CNsubsection{2.5}{REST API Design Principles}

5G SBA APIs follow RESTful conventions:

\begin{itemize}
\tightlist
\item
  \textbf{Resources} identified by URIs: \CNcode{https:/}\allowbreak{}\CNcode{/}\allowbreak{}\CNcode{\{nfInstanceId\}.}\allowbreak{}\CNcode{5gc.}\allowbreak{}\CNcode{mnc\{MNC\}.}\allowbreak{}\CNcode{mcc\{MCC\}/}\allowbreak{}\CNcode{nsmf-pdusession/}\allowbreak{}\CNcode{v1/}\allowbreak{}\CNcode{sm-contexts}
\item
  \textbf{HTTP methods} map to operations: GET (read), POST (create), PUT (update), DELETE (remove), PATCH (partial update)
\item
  \textbf{Stateless} requests — each request contains all necessary context
\item
  \textbf{HATEOAS} — responses include links to related resources
\item
  \textbf{API versioning} — version in URI path (v1, v2)
\end{itemize}

\Needspace{14\baselineskip}

\CNkeepfig{m18-sba}{8}

\CNsection{3}{5G Core Network Functions}

\CNchained\CNsubsection{}{SBA Architecture Diagram}

\CNfigure{m18-sba}{5G service-based architecture around the SBI}

\CNsubsection{3.1}{AMF — Access and Mobility Management Function}

\textbf{Purpose:} Handles all access and mobility-related signaling for the UE.

\textbf{Service Interface:} \CNcode{Namf}

\textbf{Key Services:}\CNbr{}\textbar{} Service \textbar{} Description \textbar{}\CNbr{}\textbar{}—\textbar{}—\textbar{}\CNbr{}\textbar{} \CNcode{Namf\_Communication} \textbar{} N1/\allowbreak{}N2 message transport, UE-specific signaling \textbar{}\CNbr{}\textbar{} \CNcode{Namf\_EventExposure} \textbar{} Expose mobility events (location, reachability) \textbar{}\CNbr{}\textbar{} \CNcode{Namf\_MT} \textbar{} Mobile-terminated message delivery \textbar{}\CNbr{}\textbar{} \CNcode{Namf\_Location} \textbar{} UE location services \textbar{}

\textbf{Key Responsibilities:}

\begin{itemize}
\tightlist
\item
  UE registration and deregistration
\item
  Connection management (idle \ensuremath{\leftrightarrow} connected)
\item
  Mobility management (handover, tracking area updates)
\item
  Security context management (NAS encryption/\allowbreak{}integrity)
\item
  Relay NAS-SM messages between UE and SMF
\end{itemize}

\CNsubsection{3.2}{SMF — Session Management Function}

\textbf{Purpose:} Manages PDU (data) sessions between UE and Data Network.

\textbf{Service Interface:} \CNcode{Nsmf}

\textbf{Key Services:}\CNbr{}\textbar{} Service \textbar{} Description \textbar{}\CNbr{}\textbar{}—\textbar{}—\textbar{}\CNbr{}\textbar{} \CNcode{Nsmf\_PDUSession} \textbar{} Create, modify, release PDU sessions \textbar{}\CNbr{}\textbar{} \CNcode{Nsmf\_EventExposure} \textbar{} PDU session events (e.g., IP address change) \textbar{}

\textbf{Key Responsibilities:}

\begin{itemize}
\tightlist
\item
  PDU session establishment, modification, release
\item
  UE IP address allocation and management
\item
  UPF selection and control (via N4/\allowbreak{}PFCP)
\item
  QoS flow management and enforcement
\item
  Downlink data notification
\item
  Traffic steering and routing configuration
\end{itemize}

\CNsubsection{3.3}{UPF — User Plane Function}

\textbf{Purpose:} Forwards user data packets between RAN and Data Networks.

\textbf{\CNwarn{} Note:} UPF does NOT use SBI. It uses reference points:

\begin{itemize}
\tightlist
\item
  \textbf{N3} — between RAN and UPF (GTP-U tunnel)
\item
  \textbf{N4} — between SMF and UPF (PFCP protocol)
\item
  \textbf{N6} — between UPF and Data Network (IP traffic)
\item
  \textbf{N9} — between UPFs (inter-UPF GTP-U tunnel)
\end{itemize}

\textbf{Key Responsibilities:}

\begin{itemize}
\tightlist
\item
  Packet routing and forwarding
\item
  QoS handling (marking, policing, shaping)
\item
  Traffic usage reporting
\item
  Uplink traffic verification
\item
  Lawful intercept (user plane)
\item
  Buffering of downlink packets (UE in idle mode)
\end{itemize}

\CNsubsection{3.4}{UDM — Unified Data Management}

\textbf{Purpose:} Manages subscriber data and computes authentication credentials.

\textbf{Service Interface:} \CNcode{Nudm}

\textbf{Key Services:}\CNbr{}\textbar{} Service \textbar{} Description \textbar{}\CNbr{}\textbar{}—\textbar{}—\textbar{}\CNbr{}\textbar{} \CNcode{Nudm\_}\allowbreak{}\CNcode{SubscriberDataManagement} \textbar{} Retrieve/\allowbreak{}update subscription data \textbar{}\CNbr{}\textbar{} \CNcode{Nudm\_}\allowbreak{}\CNcode{UEContextManagement} \textbar{} Track which AMF/\allowbreak{}SMF serves a UE \textbar{}\CNbr{}\textbar{} \CNcode{Nudm\_}\allowbreak{}\CNcode{UEAuthentication} \textbar{} Generate authentication vectors \textbar{}\CNbr{}\textbar{} \CNcode{Nudm\_EventExposure} \textbar{} Subscription data change notifications \textbar{}

\textbf{Key Responsibilities:}

\begin{itemize}
\tightlist
\item
  Stores subscription data (authorized slices, QoS profiles, DNN permissions)
\item
  Computes authentication credentials (using SUPI, K, OPc)
\item
  Manages UE context registration (serving AMF, serving SMF)
\item
  De-concealment of SUCI \ensuremath{\to} SUPI
\end{itemize}

\CNsubsection{3.5}{AUSF — Authentication Server Function}

\textbf{Purpose:} Executes the authentication procedure for UEs.

\textbf{Service Interface:} \CNcode{Nausf}

\textbf{Key Services:}\CNbr{}\textbar{} Service \textbar{} Description \textbar{}\CNbr{}\textbar{}—\textbar{}—\textbar{}\CNbr{}\textbar{} \CNcode{Nausf\_}\allowbreak{}\CNcode{UEAuthentication} \textbar{} Perform 5G-AKA or EAP-AKA’ authentication \textbar{}\CNbr{}\textbar{} \CNcode{Nausf\_}\allowbreak{}\CNcode{SoRProtection} \textbar{} Steering of Roaming security \textbar{}

\textbf{Key Responsibilities:}

\begin{itemize}
\tightlist
\item
  Authenticates UEs during registration
\item
  Supports 5G-AKA and EAP-AKA’ methods
\item
  Stores intermediate key material (\ensuremath{K_{\mathrm{AUSF}}})
\item
  Interfaces with UDM for authentication vectors
\end{itemize}

\CNsubsection{3.6}{PCF — Policy Control Function}

\textbf{Purpose:} Provides policy rules governing network behavior and QoS.

\textbf{Service Interface:} \CNcode{Npcf}

\textbf{Key Services:}\CNbr{}\textbar{} Service \textbar{} Description \textbar{}\CNbr{}\textbar{}—\textbar{}—\textbar{}\CNbr{}\textbar{} \CNcode{Npcf\_}\allowbreak{}\CNcode{SMPolicyControl} \textbar{} Session-level policies (QoS, charging, gating) \textbar{}\CNbr{}\textbar{} \CNcode{Npcf\_}\allowbreak{}\CNcode{AMPolicyControl} \textbar{} Access and mobility policies (RFSP, service area) \textbar{}\CNbr{}\textbar{} \CNcode{Npcf\_}\allowbreak{}\CNcode{PolicyAuthorization} \textbar{} AF-triggered policy (e.g., IMS requests QoS) \textbar{}\CNbr{}\textbar{} \CNcode{Npcf\_}\allowbreak{}\CNcode{BDTPolicyControl} \textbar{} Background data transfer policies \textbar{}\CNbr{}\textbar{} \CNcode{Npcf\_EventExposure} \textbar{} Policy-related event notifications \textbar{}

\textbf{Key Responsibilities:}

\begin{itemize}
\tightlist
\item
  QoS policy decisions
\item
  Charging policy rules
\item
  Access control and gating decisions
\item
  Network slicing policies
\item
  Usage monitoring control
\end{itemize}

\CNsubsection{3.7}{NRF — Network Repository Function}

\textbf{Purpose:} The \textbf{central registry} of all NF instances and their services — the backbone of SBA.

\textbf{Service Interface:} \CNcode{Nnrf}

\textbf{Key Services:}\CNbr{}\textbar{} Service \textbar{} Description \textbar{}\CNbr{}\textbar{}—\textbar{}—\textbar{}\CNbr{}\textbar{} \CNcode{Nnrf\_NFManagement} \textbar{} NF registration, update, deregistration \textbar{}\CNbr{}\textbar{} \CNcode{Nnrf\_NFDiscovery} \textbar{} Query available NFs by type, slice, location \textbar{}\CNbr{}\textbar{} \CNcode{Nnrf\_AccessToken} \textbar{} Issue OAuth2.0 tokens for NF-to-NF authorization \textbar{}

\begin{CNquote}

\CNstar{} \textbf{NRF is the critical enabler of SBA} — without it, NFs cannot discover each other and the service-based model collapses back to static configuration.

\end{CNquote}

\CNsubsection{3.8}{NSSF — Network Slice Selection Function}

\textbf{Purpose:} Selects the appropriate network slice instance for a UE.

\textbf{Service Interface:} \CNcode{Nnssf}

\textbf{Key Services:}\CNbr{}\textbar{} Service \textbar{} Description \textbar{}\CNbr{}\textbar{}—\textbar{}—\textbar{}\CNbr{}\textbar{} \CNcode{Nnssf\_NSSelection} \textbar{} Determine allowed slices and target AMF set \textbar{}\CNbr{}\textbar{} \CNcode{Nnssf\_}\allowbreak{}\CNcode{NSSAIAvailability} \textbar{} Manage NSSAI availability per tracking area \textbar{}

\textbf{Key Responsibilities:}

\begin{itemize}
\tightlist
\item
  Determines the Allowed NSSAI for UE registration
\item
  Selects target AMF Set/\allowbreak{}AMF based on requested slices
\item
  Manages per-TA NSSAI availability information
\end{itemize}

\CNsubsection{3.9}{NEF — Network Exposure Function}

\textbf{Purpose:} Securely exposes 5G network capabilities to external applications.

\textbf{Service Interface:} \CNcode{Nnef}

\textbf{Key Services:}\CNbr{}\textbar{} Service \textbar{} Description \textbar{}\CNbr{}\textbar{}—\textbar{}—\textbar{}\CNbr{}\textbar{} \CNcode{Nnef\_EventExposure} \textbar{} Expose network events to AFs \textbar{}\CNbr{}\textbar{} \CNcode{Nnef\_PFDManagement} \textbar{} Packet Flow Description management \textbar{}\CNbr{}\textbar{} \CNcode{Nnef\_}\allowbreak{}\CNcode{TrafficInfluence} \textbar{} Allow AF to influence traffic routing \textbar{}\CNbr{}\textbar{} \CNcode{Nnef\_}\allowbreak{}\CNcode{ParameterProvision} \textbar{} Provision expected UE behavior \textbar{}

\textbf{Key Responsibilities:}

\begin{itemize}
\tightlist
\item
  API gateway between 5G core and external world
\item
  Translates internal 5GC information to external representations
\item
  Enforces security policies on external access
\item
  Rate limiting and API management
\end{itemize}

\CNsubsection{3.10}{AF — Application Function}

\textbf{Purpose:} Represents third-party applications that interact with the 5G core.

\textbf{Service Interface:} \CNcode{Naf}

\textbf{Key Responsibilities:}

\begin{itemize}
\tightlist
\item
  Requests QoS for application flows (e.g., video streaming)
\item
  Influences traffic routing (e.g., edge computing)
\item
  Subscribes to network events (e.g., UE location)
\item
  Can access 5GC directly (trusted) or via NEF (untrusted)
\end{itemize}

\textbf{Examples:} IMS (VoNR), streaming platforms, enterprise applications, MEC orchestrators.

\CNsection{4}{NRF — The Key Enabler of SBA}

The NRF is the \textbf{most critical component} that makes SBA work. Without NRF, 5G would revert to static, pre-configured connections like EPC.

\CNsubsection{4.1}{NF Registration}

When an NF instance starts, it \textbf{registers} with the NRF by providing its \textbf{NF Profile}:

\tcbset{CNcodemode/.style={unbreakable}}\def\CNcodelang{JSON}

\begin{Shaded}
\begin{Highlighting}[]
\FunctionTok{\{}
  \DataTypeTok{"nfInstanceId"}\FunctionTok{:} \StringTok{"3fa85f64{-}5717{-}4562{-}b3fc{-}2c963f66afa6"}\FunctionTok{,}
  \DataTypeTok{"nfType"}\FunctionTok{:} \StringTok{"SMF"}\FunctionTok{,}
  \DataTypeTok{"nfStatus"}\FunctionTok{:} \StringTok{"REGISTERED"}\FunctionTok{,}
  \DataTypeTok{"heartBeatTimer"}\FunctionTok{:} \DecValTok{30}\FunctionTok{,}
  \DataTypeTok{"plmnList"}\FunctionTok{:} \OtherTok{[}\FunctionTok{\{}\DataTypeTok{"mcc"}\FunctionTok{:} \StringTok{"310"}\FunctionTok{,} \DataTypeTok{"mnc"}\FunctionTok{:} \StringTok{"014"}\FunctionTok{\}}\OtherTok{]}\FunctionTok{,}
  \DataTypeTok{"sNssais"}\FunctionTok{:} \OtherTok{[}\FunctionTok{\{}\DataTypeTok{"sst"}\FunctionTok{:} \DecValTok{1}\FunctionTok{,} \DataTypeTok{"sd"}\FunctionTok{:} \StringTok{"000001"}\FunctionTok{\}}\OtherTok{]}\FunctionTok{,}
  \DataTypeTok{"fqdn"}\FunctionTok{:} \StringTok{"smf1.5gc.mnc014.mcc310.3gppnetwork.org"}\FunctionTok{,}
  \DataTypeTok{"ipv4Addresses"}\FunctionTok{:} \OtherTok{[}\StringTok{"10.10.10.5"}\OtherTok{]}\FunctionTok{,}
  \DataTypeTok{"nfServices"}\FunctionTok{:} \OtherTok{[}
    \FunctionTok{\{}
      \DataTypeTok{"serviceInstanceId"}\FunctionTok{:} \StringTok{"nsmf{-}pdusession{-}01"}\FunctionTok{,}
      \DataTypeTok{"serviceName"}\FunctionTok{:} \StringTok{"nsmf{-}pdusession"}\FunctionTok{,}
      \DataTypeTok{"versions"}\FunctionTok{:} \OtherTok{[}\FunctionTok{\{}\DataTypeTok{"apiVersionInUri"}\FunctionTok{:} \StringTok{"v1"}\FunctionTok{,} \DataTypeTok{"apiFullVersion"}\FunctionTok{:} \StringTok{"1.2.0"}\FunctionTok{\}}\OtherTok{]}\FunctionTok{,}
      \DataTypeTok{"scheme"}\FunctionTok{:} \StringTok{"https"}\FunctionTok{,}
      \DataTypeTok{"nfServiceStatus"}\FunctionTok{:} \StringTok{"REGISTERED"}\FunctionTok{,}
      \DataTypeTok{"capacity"}\FunctionTok{:} \DecValTok{100}\FunctionTok{,}
      \DataTypeTok{"load"}\FunctionTok{:} \DecValTok{25}\FunctionTok{,}
      \DataTypeTok{"priority"}\FunctionTok{:} \DecValTok{1}
    \FunctionTok{\}}
  \OtherTok{]}\FunctionTok{,}
  \DataTypeTok{"smfInfo"}\FunctionTok{:} \FunctionTok{\{}
    \DataTypeTok{"sNssaiSmfInfoList"}\FunctionTok{:} \OtherTok{[}\FunctionTok{\{}\DataTypeTok{"sNssai"}\FunctionTok{:} \FunctionTok{\{}\DataTypeTok{"sst"}\FunctionTok{:} \DecValTok{1}\FunctionTok{\},} \DataTypeTok{"dnnSmfInfoList"}\FunctionTok{:} \OtherTok{[}\FunctionTok{\{}\DataTypeTok{"dnn"}\FunctionTok{:} \StringTok{"internet"}\FunctionTok{\}}\OtherTok{]}\FunctionTok{\}}\OtherTok{]}
  \FunctionTok{\},}
  \DataTypeTok{"capacity"}\FunctionTok{:} \DecValTok{100}\FunctionTok{,}
  \DataTypeTok{"load"}\FunctionTok{:} \DecValTok{25}\FunctionTok{,}
  \DataTypeTok{"locality"}\FunctionTok{:} \StringTok{"DC{-}East{-}1"}
\FunctionTok{\}}
\end{Highlighting}
\end{Shaded}

\textbf{NF Profile contains:}

\begin{itemize}
\tightlist
\item
  \textbf{NF Type} — what kind of NF (AMF, SMF, PCF, etc.)
\item
  \textbf{Services offered} — list of service names and versions
\item
  \textbf{Capacity} — relative capacity for load balancing (0-65535)
\item
  \textbf{Load} — current load percentage (0-100)
\item
  \textbf{Locality} — geographic/\allowbreak{}data center location for proximity-based selection
\item
  \textbf{PLMN} — which networks this NF serves
\item
  \textbf{S-NSSAI} — which slices this NF supports
\item
  \textbf{IP/\allowbreak{}FQDN} — how to reach this NF
\end{itemize}

\CNsubsection{4.2}{NF Discovery}

When an NF needs a service, it queries the NRF:

\tcbset{CNcodemode/.style={unbreakable}}\def\CNcodelang{}

\begin{CNplaincode}
\begin{Verbatim}
GET /nnrf-disc/v1/nf-instances?target-nf-type=SMF&requester-nf-type=AMF&snssai={"sst":1,"sd":"000001"}&dnn=internet
\end{Verbatim}
\end{CNplaincode}

\textbf{Discovery query parameters:}\CNbr{}\textbar{} Parameter \textbar{} Purpose \textbar{}\CNbr{}\textbar{}—\textbar{}—\textbar{}\CNbr{}\textbar{} \CNcode{target-nf-type} \textbar{} Type of NF being searched (mandatory) \textbar{}\CNbr{}\textbar{} \CNcode{requester-nf-type} \textbar{} Type of the requesting NF (mandatory) \textbar{}\CNbr{}\textbar{} \CNcode{service-names} \textbar{} Specific services needed \textbar{}\CNbr{}\textbar{} \CNcode{snssai} \textbar{} Network slice filter \textbar{}\CNbr{}\textbar{} \CNcode{dnn} \textbar{} Data Network Name filter \textbar{}\CNbr{}\textbar{} \CNcode{tai} \textbar{} Tracking area for locality \textbar{}\CNbr{}\textbar{} \CNcode{preferred-locality} \textbar{} Prefer NFs in same data center \textbar{}

\textbf{NRF returns:} Ranked list of matching NF instances with endpoints and load information.

\CNsubsection{4.3}{NF Status Notification}

NFs can \textbf{subscribe} to NRF for status events:

\begin{itemize}
\tightlist
\item
  NF registered /\allowbreak{} deregistered
\item
  NF profile changed (load, capacity)
\item
  NF service status changed
\end{itemize}

This enables:

\begin{itemize}
\tightlist
\item
  \textbf{AMF} to be notified when new SMFs become available
\item
  \textbf{Load balancers} to update routing when NF load changes
\item
  \textbf{Fault detection} when NF heartbeat expires
\end{itemize}

\CNkeepfig{m18-nrf-discovery}{3}

\CNsubsection{}{NF Discovery Flow Diagram}

\CNfigure{m18-nrf-discovery}{NF registration, discovery and subscription through the NRF}

\Needspace{14\baselineskip}

\CNsection{5}{Service Discovery and Communication Flow}

\CNchained\CNsubsection{5.1}{Direct Communication}

In the \textbf{direct model}, NF-A communicates directly with NF-B after discovery:

\tcbset{CNcodemode/.style={unbreakable}}\def\CNcodelang{}

\begin{CNplaincode}
\begin{Verbatim}
1. NF-A → NRF: "Find me an SMF for slice 1, DNN internet"
2. NRF → NF-A: "Here are 3 SMFs: SMF-1 (load=20%), SMF-2 (load=45%), SMF-3 (load=80%)"
3. NF-A selects SMF-1 (lowest load)
4. NF-A → SMF-1: Direct service request
\end{Verbatim}
\end{CNplaincode}

\textbf{Advantages:} Low latency (no intermediate hop), simple\CNbr{}\textbf{Disadvantages:} NF-A must handle load balancing, retries, failover logic

\CNkeepfig{m18-scp}{5}

\CNsubsection{5.2}{Indirect Communication via SCP}

The \textbf{Service Communication Proxy (SCP)} acts as an intermediary:

\CNfigure{m18-scp}{Indirect communication through the SCP}

\textbf{SCP provides:}

\begin{itemize}
\tightlist
\item
  Delegated discovery (SCP queries NRF on behalf of NF-A)
\item
  Load balancing across NF-B instances
\item
  Message routing and topology hiding
\item
  Retry and failover logic
\item
  Monitoring and observability
\item
  Connection pooling
\end{itemize}

\textbf{Communication Models (3GPP TS 23.501):}\CNbr{}\textbar{} Model \textbar{} Discovery \textbar{} Routing \textbar{}\CNbr{}\textbar{}—\textbar{}—\textbar{}—\textbar{}\CNbr{}\textbar{} Model A \textbar{} NF-A via NRF \textbar{} Direct NF-A \ensuremath{\to} NF-B \textbar{}\CNbr{}\textbar{} Model B \textbar{} NF-A via NRF \textbar{} Via SCP (NF-A \ensuremath{\to} SCP \ensuremath{\to} NF-B) \textbar{}\CNbr{}\textbar{} Model C \textbar{} Delegated to SCP \textbar{} Via SCP (NF-A \ensuremath{\to} SCP \ensuremath{\to} NF-B) \textbar{}\CNbr{}\textbar{} Model D \textbar{} Delegated to SCP \textbar{} Direct NF-A \ensuremath{\to} NF-B \textbar{}

\Needspace{14\baselineskip}

\CNkeepfig{m18-sbi-stack}{8}

\CNsection{6}{SBI Protocol Stack}

\CNchained\CNsubsection{}{Complete Protocol Stack}

\CNfigure{m18-sbi-stack}{SBI protocol stack}

\CNsubsection{}{OAuth 2.0 for NF Authorization}

Before NF-A can call NF-B’s service, it must obtain an \textbf{access token} from the NRF:

\tcbset{CNcodemode/.style={unbreakable}}\def\CNcodelang{}

\begin{CNplaincode}
\begin{Verbatim}
1. NF-A → NRF: POST /oauth2/token
   {
     "grant_type": "client_credentials",
     "nfInstanceId": "{NF-A ID}",
     "nfType": "AMF",
     "targetNfType": "SMF",
     "scope": "nsmf-pdusession"
   }
   
2. NRF → NF-A: Access Token (JWT)
   {
     "access_token": "eyJ...",
     "token_type": "Bearer",
     "expires_in": 3600,
     "scope": "nsmf-pdusession"
   }

3. NF-A → NF-B: Service Request + Authorization: Bearer eyJ...
4. NF-B validates the token (signature, expiry, scope)
\end{Verbatim}
\end{CNplaincode}

\Needspace{14\baselineskip}

\CNkeepfig{m18-epc-vs-sba}{8}

\CNsection{7}{Comparison: EPC Interfaces vs 5G SBA}

\CNchained\CNsubsection{}{EPC vs SBA Comparison Diagram}

\CNfigure{m18-epc-vs-sba}{Point-to-point EPC interfaces versus the SBA bus}

\Needspace{26\baselineskip}

\CNsubsection{}{Detailed Comparison Table}

{\def\LTcaptype{none} % do not increment counter
\Needspace{22\baselineskip}
\begin{longtable}[c]{@{}
  >{\raggedright\arraybackslash\hspace{0pt}}m{(\linewidth - 4\tabcolsep) * \real{0.3333}}
  >{\raggedright\arraybackslash\hspace{0pt}}m{(\linewidth - 4\tabcolsep) * \real{0.3333}}
  >{\raggedright\arraybackslash\hspace{0pt}}m{(\linewidth - 4\tabcolsep) * \real{0.3333}}@{}}
\toprule\noalign{}
\rowcolor{CNink}\CNth{Aspect} & \CNth{4G EPC} & \CNth{5G SBA} \\
\CNheadrulesep
\endhead
\bottomrule\noalign{}
\endlastfoot
\textbf{Architecture} & Point-to-point reference points & Service-based interfaces on common bus \\
\textbf{Protocols} & Diameter, GTP-C, S1AP (multiple) & HTTP/\allowbreak{}2 + JSON (unified) \\
\textbf{Coupling} & Tight — each pair has unique interface & Loose — common framework for all \\
\textbf{Adding new NF} & Define new interfaces to all peers & Register services with NRF \\
\textbf{Discovery} & Static configuration & Dynamic via NRF \\
\textbf{Scaling} & Scale entire monolithic element & Scale individual NF services independently \\
\textbf{Deployment} & Hardware appliances /\allowbreak{} VMs & Cloud-native containers (Kubernetes) \\
\textbf{State} & Stateful network elements & Stateless NFs (externalized state) \\
\textbf{Redundancy} & Active-standby pairs & Horizontal scaling + NRF re-selection \\
\textbf{Load balancing} & DNS-based or static & NRF-based (capacity, load, locality) \\
\textbf{Authorization} & Hop-by-hop TLS & End-to-end OAuth 2.0 tokens \\
\textbf{API definition} & 3GPP-specific ASN.1/\allowbreak{}XML & OpenAPI 3.0 (industry standard) \\
\textbf{Extensibility} & New release = new interfaces & New services registered dynamically \\
\textbf{Vendor ecosystem} & Telecom-specific vendors & Cloud/\allowbreak{}IT vendors can participate \\
\end{longtable}
}

\Needspace{20\baselineskip}

\CNsubsection{}{Functional Mapping: EPC \ensuremath{\to} 5G}

{\def\LTcaptype{none} % do not increment counter
\Needspace{16\baselineskip}
\begin{longtable}[c]{@{}ccc@{}}
\toprule\noalign{}
\rowcolor{CNink}\CNth{EPC Element} & \CNth{5G Equivalent} & \CNth{Key Difference} \\
\CNheadrulesep
\endhead
\bottomrule\noalign{}
\endlastfoot
MME & AMF + SMF & Separated access/\allowbreak{}mobility from session mgmt \\
S-GW + P-GW & UPF & Simplified to single user plane element \\
HSS & UDM + AUSF & Separated data storage from authentication \\
PCRF & PCF & Same role, now on SBI \\
— (new) & NRF & Service registry (no EPC equivalent!) \\
— (new) & NSSF & Slice selection (no slicing in 4G) \\
— (new) & NEF & Replaces SCEF with richer API exposure \\
\end{longtable}
}

\Needspace{20\baselineskip}

\CNsection{}{Summary}

{\def\LTcaptype{none} % do not increment counter
\Needspace{14\baselineskip}
\begin{longtable}[c]{@{}
  >{\raggedright\arraybackslash\hspace{0pt}}m{(\linewidth - 2\tabcolsep) * \real{0.5000}}
  >{\raggedright\arraybackslash\hspace{0pt}}m{(\linewidth - 2\tabcolsep) * \real{0.5000}}@{}}
\toprule\noalign{}
\rowcolor{CNink}\CNth{Concept} & \CNth{Key Takeaway} \\
\CNheadrulesep
\endhead
\bottomrule\noalign{}
\endlastfoot
\textbf{SBA motivation} & EPC’s N\ensuremath{\times}(N-1)/\allowbreak{}2 interfaces don’t scale; SBA uses a common bus \\
\textbf{Design principles} & HTTP/\allowbreak{}2, JSON, REST, microservices, stateless \\
\textbf{NRF} & Heart of SBA — registration, discovery, authorization \\
\textbf{Communication} & Direct (NF\ensuremath{\to}NF) or Indirect (via SCP) \\
\textbf{Security} & mTLS + OAuth 2.0 tokens from NRF \\
\textbf{Key difference} & EPC = static, rigid, monolithic; SBA = dynamic, flexible, cloud-native \\
\end{longtable}
}

\CNboxsection{Review Questions}

\begin{CNbox}[unbreakable]{q}{Review Questions}

\begin{enumerate}
\def\labelenumi{\arabic{enumi}.}
\tightlist
\item
  Calculate how many unique interfaces are needed for 8 EPC elements. How does SBA solve this?
\item
  Why was HTTP/\allowbreak{}2 chosen over HTTP/\allowbreak{}1.1 for SBI?
\item
  What information is included in an NF Profile during NRF registration?
\item
  Why is the UPF not connected to the SBI bus?
\item
  Compare SCP communication Models B and C — when would you use each?
\item
  How does OAuth 2.0 work in the 5G SBA context? What role does NRF play?
\item
  An operator wants to add a new Analytics NF (NWDAF). Compare the effort in EPC vs SBA architecture.
\item
  Why is NRF considered the ``single most critical'' component of SBA?
\end{enumerate}

\end{CNbox}

\CNsection{}{References}

\begin{itemize}
\tightlist
\item
  3GPP TS 23.501 — System Architecture for the 5G System (Stage 2)
\item
  3GPP TS 23.502 — Procedures for the 5G System (Stage 2)
\item
  3GPP TS 29.510 — NRF Services (Stage 3)
\item
  3GPP TS 29.500 — Technical Realization of Service-Based Architecture
\item
  3GPP TS 33.501 — Security Architecture and Procedures for 5G System
\end{itemize}
